\documentclass[12pt,a4paper]{article}

\usepackage[utf8]{inputenc}
\usepackage[T1]{fontenc}
\usepackage[provide=*,spanish]{babel}
\usepackage{mathpazo}
\usepackage{microtype}
\usepackage{amsmath,amssymb,amsthm}
\usepackage{booktabs}
\usepackage{tabularx}
\usepackage{geometry}
\usepackage{enumitem}
\usepackage[hidelinks]{hyperref}

\geometry{margin=2.5cm}

\title{Ecuaciones Diferenciales Ordinarias\\Clase III -- Teórico}
\author{Benjamín Marcolongo}
\date{}

\begin{document}

	\maketitle

	\section*{Introducción}

	En la Clase II estudiamos qué información puede obtenerse de una ecuación diferencial antes de resolverla: dominio, existencia y unicidad, formas de representación y campos de direcciones. En esta clase comenzaremos a desarrollar métodos analíticos para construir soluciones.

	En todo el material usaremos \(\dot y\) cuando la variable independiente sea el tiempo. Si la variable independiente es \(x\), escribiremos explícitamente \(dy/dx\).

	\clearpage
	\section{Resolución por integración directa}

	El caso más sencillo de una EDO de primer orden es

	\[
	\frac{dy}{dx}=g(x),
	\]
	%
	donde la derivada depende únicamente de la variable independiente. Si \(G\) es una primitiva de \(g\), entonces

	\[
	y(x)=G(x)+C.
	\]
	%
	La constante \(C\) representa una familia uniparamétrica de soluciones. Una condición inicial permite seleccionar un único miembro de esa familia.

	\subsection{Problema de valor inicial}

	Consideremos

	\[
	\begin{cases}
	\dfrac{dy}{dx}=x^2+1,\\[2mm]
	y(0)=1.
	\end{cases}
	\]
	%
	Integrando,

	\[
	y(x)=\frac{x^3}{3}+x+C.
	\]
	%
	La condición inicial implica

	\[
	1=y(0)=C.
	\]
	%
	Por lo tanto,

	\[
	\boxed{y(x)=\frac{x^3}{3}+x+1.}
	\]

	También puede incorporarse directamente la condición inicial mediante una integral definida. Para

	\[
	\begin{cases}
	\dfrac{dy}{dx}=g(x),\\[2mm]
	y(x_0)=y_0,
	\end{cases}
	\]
	%
	se obtiene

	\[
	\boxed{y(x)=y_0+\int_{x_0}^{x}g(s)\,ds.}
	\]
	%
	Esta expresión deja explícita la dependencia de la solución respecto de la condición inicial.

	\subsection{El intervalo de definición}

	La integración no elimina las restricciones del dominio. Por ejemplo,

	\[
	\frac{dy}{dx}=\frac{1}{x}
	\]
	%
	tiene como familia de soluciones

	\[
	y(x)=\ln|x|+C,
	\]
	%
	pero ninguna solución puede estar definida en un intervalo que atraviese \(x=0\). Si se impone una condición inicial con \(x_0<0\), el intervalo maximal correspondiente está contenido en \(( -\infty,0 )\); si \(x_0>0\), está contenido en \((0,\infty)\).

	\clearpage
	\section{Ecuaciones de variables separables}

	Una ecuación diferencial de primer orden es \textbf{separable} si puede escribirse como

	\[
	\frac{dy}{dx}=g(x)h(y).
	\]
	%
	Antes de dividir por \(h(y)\), deben buscarse sus ceros. Si

	\[
	h(y_\ast)=0,
	\]
	%
	entonces \(y(x)=y_\ast\) es una solución de equilibrio. Dividir por \(h(y)\) elimina esa posibilidad y puede hacer que se pierdan soluciones.

	En un intervalo donde \(h(y)\neq0\), podemos escribir

	\[
	\frac{1}{h(y)}\,dy=g(x)\,dx.
	\]

	\subsection{Justificación de la separación de variables}

	La expresión anterior parece obtenida pasando \(dx\) y \(dy\) de un miembro al otro como si fueran factores algebraicos. Sin embargo, \(dy/dx\) no necesita interpretarse como una fracción ordinaria. El procedimiento se justifica mediante la regla de la cadena.

	Sea \(y=y(x)\) una solución y supongamos que \(h(y(x))\neq0\) en el intervalo considerado. Al dividir la ecuación por \(h(y(x))\), obtenemos

	\[
	\frac{1}{h(y(x))}\frac{dy}{dx}=g(x).
	\]
	%
	Definamos una función \(H\) tal que

	\[
	H'(y)=\frac{1}{h(y)}.
	\]
	%
	Entonces, por la regla de la cadena,

	\[
	\frac{d}{dx}H(y(x))
	=
	H'(y(x))\frac{dy}{dx}
	=
	\frac{1}{h(y(x))}\frac{dy}{dx}.
	\]
	%
	Por lo tanto, la ecuación separable equivale a

	\[
	\frac{d}{dx}H(y(x))=g(x).
	\]
	%
	Integrando respecto de \(x\),

	\[
	H(y(x))=G(x)+C,
	\qquad
	G'(x)=g(x).
	\]
	%
	La notación

	\[
	\frac{dy}{h(y)}=g(x)\,dx
	\]
	%
	es una forma abreviada de este argumento. Sobre la curva solución se cumple \(dy=(dy/dx)\,dx\), y la igualdad entre diferenciales expresa precisamente la regla de la cadena.

	\begin{quote}
		\textbf{Importante.} Separar variables no consiste en tratar \(dy\) y \(dx\) como números independientes. Es una notación compacta para dividir por \(h(y)\), aplicar la regla de la cadena e integrar. La división solo es válida donde \(h(y)\neq0\); por eso los equilibrios deben estudiarse antes.
	\end{quote}

	Integrando ambos miembros,

	\[
	\int\frac{1}{h(y)}\,dy
	=
	\int g(x)\,dx+C.
	\]
	%
	El resultado puede definir a \(y\) explícitamente o mediante una relación implícita. Despejar \(y\) es conveniente cuando puede hacerse, pero no es necesario para caracterizar una curva solución.

	\subsection{Procedimiento}

	Para resolver una ecuación separable conviene seguir este orden:

	\begin{enumerate}[label=\arabic*.]
		\item escribirla en la forma \(dy/dx=g(x)h(y)\);
		\item encontrar las soluciones de equilibrio resolviendo \(h(y)=0\);
		\item separar las variables en los intervalos donde \(h(y)\neq0\);
		\item integrar ambos miembros;
		\item despejar \(y\), si es posible;
		\item imponer la condición inicial;
		\item determinar el intervalo de definición de la solución obtenida.
	\end{enumerate}

	\clearpage
	\subsection{Ejemplo: una solución que puede perderse}

	Consideremos

	\[
	\frac{dy}{dx}=3y+1.
	\]
	%
	La ecuación es autónoma y separable. Antes de dividir, observamos que

	\[
	3y+1=0
	\quad\Longrightarrow\quad
	y=-\frac13
	\]
	%
	es una solución de equilibrio. Para las soluciones no constantes,

	\[
	\frac{dy}{3y+1}=dx.
	\]
	%
	Integrando,

	\[
	\frac13\ln|3y+1|=x+C.
	\]
	%
	Equivalentemente,

	\[
	3y+1=Ke^{3x},
	\]
	%
	por lo que

	\[
	\boxed{y(x)=\frac{Ke^{3x}-1}{3}.}
	\]
	%
	En este caso, \(K=0\) recupera la solución de equilibrio. Esto no ocurre siempre: por eso los equilibrios deben buscarse antes de dividir.

	Si además se impone \(y(0)=1\), entonces \(K=4\) y

	\[
	y(x)=\frac{4e^{3x}-1}{3}.
	\]

	\subsection{Ejemplo: dominio e intervalo maximal}

	Consideremos

	\[
	\frac{dy}{dx}=2xy^2.
	\]
	%
	La solución de equilibrio es \(y=0\). Para \(y\neq0\),

	\[
	\frac{dy}{y^2}=2x\,dx.
	\]
	%
	Integrando,

	\[
	-\frac{1}{y}=x^2+C,
	\]
	%
	y, por lo tanto,

	\[
	\boxed{y(x)=-\frac{1}{x^2+C}.}
	\]
	%
	La expresión solo define una solución en intervalos que no contengan ceros del denominador. La fórmula algebraica, por sí sola, no especifica el intervalo de definición: ese intervalo forma parte de la solución.

	\clearpage
	\section{Ecuaciones autónomas}

	Una ecuación diferencial de primer orden es \textbf{autónoma} si la variable independiente no aparece explícitamente. Su forma general es

	\[
	\dot y=f(y).
	\]
	%
	Si la variable independiente es el tiempo, la tasa de cambio depende únicamente del estado actual del sistema y no del instante en el que se lo observa.

	Por ejemplo,

	\[
	\dot y=y(1-y)
	\]
	%
	es autónoma, mientras que

	\[
	\dot y=t-y
	\]
	%
	no lo es, porque el tiempo aparece explícitamente.

	\subsection{Soluciones de equilibrio}

	Un valor \(y_\ast\) es un \textbf{punto de equilibrio} si

	\[
	f(y_\ast)=0.
	\]
	%
	En ese caso, la función constante

	\[
	y(t)=y_\ast
	\]
	%
	es solución de la ecuación. Por ejemplo, para

	\[
	\dot y=y(1-y),
	\]
	%
	los valores \(y_\ast=0\) e \(y_\ast=1\) determinan las soluciones constantes

	\[
	y(t)=0,
	\qquad
	y(t)=1.
	\]

	La clasificación de estos equilibrios mediante estabilidad se desarrollará más adelante. Por ahora, lo importante es identificarlos antes de realizar cualquier división algebraica, porque de otro modo pueden perderse soluciones.

	\clearpage
	\section*{Articulación con la Guía 1}

	Los problemas 18--20 practican integración directa, separación de variables y problemas de valor inicial. La definición de ecuación autónoma prepara el análisis de los problemas 15 y 16; la clasificación de los puntos de equilibrio se incorporará cuando estudiemos estabilidad.

	\subsection*{Nota sobre la elaboración del material}

	Este material fue elaborado por Benjamín Marcolongo con la asistencia de \textbf{GPT-5.6 Sol}, un modelo de OpenAI utilizado a través del entorno Codex, para la organización, redacción y revisión del contenido.

\end{document}
